\subsection{CLOSED SETS OF ROOTS}

DEFINITION 4. Let P be a subset of R.

\begin{enumerate}
\def\labelenumi{(\roman{enumi})}
\item
  P is said to be closed if the conditions \(\alpha \in \mathrm{P},\beta \in \mathrm{P},\alpha +\beta \in \mathrm{R}\) imply \(\alpha +\beta \in \mathrm{P}\) .
\item
  P is said to be parabolic if P is closed and if \(\mathrm{P} \cup (-\mathrm{P}) = \mathrm{R}\) .
\item
  P is said to be symmetric if \(\mathrm{P} = -\mathrm{P}\) .
\end{enumerate}

Lemma 3. Let C be a chamber of R and P a closed subset of R containing \(\mathrm{R}_{+}(\mathrm{C})\) (in the notation of no. 6). Let \(\Sigma = \mathrm{B}(\mathrm{C}) \cap (-\mathrm{P})\) , and let Q be the set of roots that are linear combinations of elements of \(\Sigma\) with non-positive integer coefficient. Then, \(\mathrm{P} = \mathrm{R}_{+}(\mathrm{C}) \cup \mathrm{Q}\) .

It is enough to show that \(P \cap (-R_{+}(C)) = Q\) . Let \(-\alpha \in Q\) . Then \(\alpha\) is the sum of n elements of \(\Sigma\) . We show, by induction on n, that \(-\alpha \in P\) . This is clear if n = 1. If n \textgreater{} 1, then by Prop. 19 of no. 6 we can write \(\alpha = \beta + \gamma\) with \(\gamma \in \Sigma\) and \(\beta\) the sum of n - 1 elements of \(\Sigma\) . By the induction hypothesis, \(-\beta \in P\) ; since \(-\gamma \in P\) and since P is closed, \(-\alpha \in P\) . Thus, \(Q \subseteq P \cap (-R_{+}(C))\) . Conversely, let \(-\alpha \in P \cap (-R_{+}(C))\) . Then \(\alpha\) is the sum of p elements of B(C). We show, by induction on p, that \(-\alpha \in Q\) . This is clear if p = 1. If p \textgreater{} 1, then by Prop. 19, we can write \(\alpha = \beta + \gamma\) with \(\gamma \in B(C)\) and \(\beta\) a root that is the sum of p - 1 elements of B(C). Since \(-\gamma = \beta + (-\alpha)\) and since P is closed, \(-\gamma \in P\) , hence \(\gamma \in \Sigma\) . Moreover, \(-\beta = \gamma + (-\alpha) \text{ so } -\beta \in \mathbf{P} \text{ since P is closed. By the induction hypothesis, } -\beta \in \mathbf{Q}, \text{ so } -\alpha = -\beta - \gamma \in \mathbf{Q}\) . Thus, \(\mathrm{P} \cap (-\mathrm{R}_{+}(\mathrm{C})) \subseteq \mathrm{Q}\) .

PROPOSITION 20. Let P be a subset of R. The following conditions are equivalent:

\begin{enumerate}
\def\labelenumi{(\roman{enumi})}
\item
  P is parabolic;
\item
  P is closed and there exists a chamber C of R such that \(\mathrm{P} \supseteq \mathrm{R}_{+}(\mathrm{C})\) ;
\item
  there exist a chamber C of R and a subset \(\Sigma\) of B(C) such that P is the union of \(\mathrm{R}_{+}(\mathrm{C})\) and the set Q of roots that are linear combinations of elements of \(\Sigma\) with non-positive integer coefficients.
\item
  \(\Longrightarrow\) (iii): this follows from Lemma 3.
\item
  \(\Longrightarrow\) (i): we adopt the assumptions and notation of (iii). It is clear that \(\mathrm{P} \cup (-\mathrm{P}) = \mathrm{R}\) . We show that, if \(\alpha, \beta \in \mathrm{P}\) are such that \(\alpha + \beta \in \mathrm{R}\) , then \(\alpha + \beta \in \mathrm{P}\) . This is obvious if the root \(\alpha + \beta\) is positive. Assume that \(\alpha + \beta\) is negative. Then \(\alpha + \beta = \sum_{\gamma \in \mathrm{B}(\mathbb{C})} n_{\gamma} \gamma\) , with \(n_{\gamma} \leqslant 0\) . But the coefficient of every element \(\gamma\) of \(\mathrm{B}(\mathbb{C}) - \Sigma\) in \(\alpha\) or \(\beta\) is \(\geqslant 0\) ; hence \(n_{\gamma} = 0\) if \(\gamma \in \mathrm{B}(\mathbb{C}) - \Sigma\) , so \(\alpha + \beta \in \mathbb{Q} \subseteq \mathbb{P}\) .
\item
  \(\Longrightarrow\) (ii): assume that P is parabolic. Let C be a chamber such that \(\operatorname{Card}(P\cap R_{+}(C))\) is as large as possible. Let \(\alpha\in\operatorname{B}(C)\) and assume that \(\alpha\in P\) , so that \(-\alpha\in P\) . For all \(\beta\in P\cap R_{+}(C)\) , \(\beta\) is not proportional to \(\alpha\) (for the hypothesis \(\beta=2\alpha\) would imply that \(\alpha=2\alpha+(-\alpha)\in P\) since P is closed). Thus \(s_{\alpha}(\beta)\in\operatorname{R}_{+}(C)\) (no. 6, Cor. 1 of Prop. 17). If we put \(C'=s_{\alpha}(C)\) , then \(\beta=s_{\alpha}(s_{\alpha}(\beta))\in s_{\alpha}(\operatorname{R}_{+}(C))=\operatorname{R}_{+}(C)\) , so \(-\alpha\in P\cap R_{+}(C')\) and hence \(\operatorname{Card}(P\cap R_{+}(C'))>\operatorname{Card}(P\cap R_{+}(C))\) . This is absurd, since \(\alpha\in P\) . Thus \(B(C)\subseteq P\) , and consequently \(R_{+}(C)\subseteq P\) by Prop. 19 and the fact that P is closed.
\end{enumerate}

COROLLARY 1. Let P be a subset of R. The following conditions are equivalent:

\begin{enumerate}
\def\labelenumi{(\roman{enumi})}
\item
  there exists a chamber C such that \(\mathrm{P} = \mathrm{R}_{+}(\mathrm{C})\) ;
\item
  P is closed and \(\{P, -P\}\) is a partition of R.
\end{enumerate}

The chamber C such that \(\mathrm{P} = \mathrm{R}_{+}(\mathrm{C})\) is then unique.

If \(\mathrm{P} = \mathrm{R}_{+}(\mathrm{C})\) , \(\mathrm{C}^{-}\) is the set of \(x^{*} \in \mathrm{V}^{*}\) such that \(\langle x^{*}, x \rangle > 0\) for all \(x \in \mathrm{P}\) , hence the uniqueness of \(\mathbf{C}\) .

COROLLARY 2. Assume that V is equipped with the structure of an ordered vector space such that, for this structure, every root of R is either positive or negative. Let P be the set of positive roots for this structure. Then there exists a unique chamber C of R such that \(P = R_{+}(C)\) .

Indeed, P satisfies condition (ii) of Cor. 1.

This corollary applies in particular when the order being considered is total, the condition on R then being automatically satisfied. Recall that such an order can be obtained, for example, by choosing a basis \((e_{i})_{1\leqslant i\leqslant n}\) of V and taking the lexicographic order on V, so \(x=\sum_{i}\xi_{i}e_{i}\) is \(\geqslant0\) if all the \(\xi_{i}\) are 0, or if \(\xi_{i}>0\) for the smallest index i such that \(\xi_{i}\neq0\) .

COROLLARY 3. A subset B of R is a basis of R if and only if the following conditions are satisfied:

\begin{enumerate}
\def\labelenumi{(\roman{enumi})}
\item
  the elements of B are linearly independent;
\item
  every root of R is a linear combination of elements of B in which the coefficients are either all positive or all negative;
\item
  every root of B is indivisible.
\end{enumerate}

We already know that the conditions are necessary (no. 5, Th. 2, and no. 6, Th. 3). Assume that conditions (i), (ii), (iii) are satisfied. Let P be the set of roots that are linear combinations of elements of B with coefficients \(\geqslant0\) . Since P satisfies condition (ii) of Cor. 1, there exists a chamber C such that \(P = R_{+}(C)\) ; let \(B' = B(C)\) , and let X and \(X'\) be the convex cones generated by B and \(B'\) . Then

\[
\mathrm{B} \subseteq \mathrm{P} \subseteq \mathrm{X} \quad \text { and } \quad \mathrm{B} ^ {\prime} \subseteq \mathrm{P} \subseteq \mathrm{X} ^ {\prime},
\]

which shows that X and \(X'\) are both generated by P, and hence coincide. But the half-lines generated by the elements of B (resp. by \(B'\) ) are the extreme generators of X (resp. \(X'\) ); since such a half-line contains only one indivisible root, \(B = B'\) .

COROLLARY 4. Let B be a basis of R, B' a subset of B, V' the vector subspace of V generated by B', and R' = R ∩ V'. Then B' is a basis of the root system R'.

This follows immediately from Cor. 3 and the Cor. of Prop. 4.

We call R' the root system generated by B'.

COROLLARY 5. Let B be a basis of R, \(A_{1}, A_{2}, \ldots, A_{r}\) pairwise orthogonal subsets of B, and \(A = A_{1} \cup A_{2} \cup \cdots \cup A_{r}\) . Then every root \(\alpha\) that is a linear combination of elements of A is actually a linear combination of elements of one of the \(A_{i}\) . In particular, if R is irreducible, there is no partition of B into pairwise orthogonal subsets.

Let \(E_1, \ldots, E_r\) , \(E\) be the vector subspaces of \(V\) generated by \(A_1, \ldots, A_r\) , \(A\) , respectively. By Cor. 4, we can assume that \(E = V\) . Then, by Th. 2 (vii) of no. 5, the \(E_i\) are stable under \(W(R)\) , so \(R\) is the union of the \(R \cap E_i\) (no. 2, Prop. 5).

COROLLARY 6. We adopt the hypotheses and notation of Prop. 20. Let \(V_{1}\) be the vector subspace of V generated by \(\Sigma\) . Then \(P \cap (-P) = Q \cup (-Q) = V_{1} \cap R\) is a root system in \(V_{1}\) with basis \(\Sigma\) .

We have \(P \cap (-P) = (R_{+}(C) \cup Q) \cap ((-R_{+}(C)) \cup (-Q)) = Q \cup (-Q)\) . Th. 3 proves that \(Q \cup (-Q) = V_{1} \cap R\) . Finally, \(\Sigma\) is a basis of the root system \(V_{1} \cap R\) by Cor. 4.

PROPOSITION 21. Let C (resp. C') be a chamber of R, \(\Sigma\) (resp. \(\Sigma'\) ) a subset of B(C) (resp. B(C')), Q (resp. Q') the set of linear combinations of elements of \(\Sigma\) (resp. \(\Sigma'\) ) with negative integer coefficients, and P = Q ∪ R \(_{+}\) (C) (resp. P' = Q' ∪ R \(_{+}\) (C')). If there exists an element of the Weyl group transforming P to P', then there exists an element of the Weyl group transforming C to C' and \(\Sigma\) to \(\Sigma'\) .

We reduce immediately to the case \(P = P'\) . Let \(V_1\) be the vector subspace of \(V\) generated by \(P \cap (-P)\) . Then \(\Sigma\) and \(\Sigma'\) are bases of the root system \(R_1 = P \cap (-P)\) in \(V_1\) (Cor. 6 of Prop. 20). Hence there exists \(g_1 \in W(R_1)\) such that \(g_1(\Sigma) = \Sigma'\) . It is clear that \(g_1\) is induced by an element \(g\) of \(W(R)\) , that is a product of the symmetries \(s_\sigma\) with \(\sigma \in \Sigma\) . Let \(\gamma = \sum_{\beta \in B(C)} c_\beta \beta\) be an element of \(P - R_1\) . Then \(c_\beta > 0\) for at least one \(\beta \in B(C) - \Sigma\) . Moreover, if \(\sigma \in \Sigma\) , then \(s_\sigma(\gamma) - \gamma \in V_1\) , so \(s_\sigma(\gamma)\) has at least one coordinate \(> 0\) with respect to \(B(C)\) (no. 1, formula (5)), hence \(s_\sigma(\gamma) \in R_+(C)\) and finally \(s_\sigma(\gamma) \in P - R_1\) . It follows that \(P - R_1\) is stable under the \(s_\sigma, \sigma \in \Sigma\) , and hence under \(g\) , so \(g(P) = P\) . We are thus reduced to proving the proposition when \(P = P'\) and \(\Sigma = \Sigma'\) . In this case, \(Q = Q'\) , so \(R_+(C) = P - Q = P - Q' = R_+(C')\) , and hence \(C = C'\) (Cor. 1 of Prop. 20).

COROLLARY. Let \(\mathrm{P}, \mathrm{P}'\) be two parabolic subsets of \(\mathbb{R}\) transformed into each other by an element of the Weyl group. If there exists a chamber \(\mathbb{C}\) of \(\mathbb{R}\) such that \(\mathrm{R}_{+}(\mathbb{C}) \subseteq \mathrm{P}\) and \(\mathrm{R}_{+}(\mathbb{C}) \subseteq \mathrm{P}'\) , then \(\mathrm{P} = \mathrm{P}'\) .

This follows from Lemma 3 and Prop. 21 since the only element of W(R) transforming C to C is 1, cf.~no. 5, Th. 2.

PROPOSITION 22. Let P be a closed subset of R such that \(P \cap (-P) = \varnothing\) . Then there exists a chamber C of R such that \(P \subseteq R_{+}(C)\) .

\begin{enumerate}
\def\labelenumi{\arabic{enumi})}
\item
  In view of the Cor. of Th. 1, no. 3, the assumptions \(\alpha \in \mathrm{P}\) , \(\beta \in \mathrm{P}\) , \((\alpha |\beta) < 0\) imply that \(\alpha + \beta \in \mathrm{P}\) .
\item
  We show that no sum \(\alpha_{1} + \cdots + \alpha_{q}\) ( \(q \geqslant 1\) ) of elements of P is zero. We proceed by induction on q. The assertion being clear for q = 1, assume that \(q \geqslant 2\) . If \(\alpha_{1} + \cdots + \alpha_{q} = 0\) , then
\end{enumerate}

\[
\alpha_ {1} = \alpha_ {2} + \dots + \alpha_ {q},
\]

so \((-\alpha_{1}|\alpha_{2}+\cdots+\alpha_{q})>0\) , hence there exists \(j\in[2,q]\) such that \((\alpha_{1}|\alpha_{j})<0\) . By part 1) of the proof, \(\alpha_{1}+\alpha_{j}\in P\) , and the relation \((\alpha_{1}+\alpha_{j})+\sum_{i\neq1,j}\alpha_{i}=0\) contradicts the induction hypothesis.

\begin{enumerate}
\def\labelenumi{\arabic{enumi})}
\setcounter{enumi}{2}
\tightlist
\item
  We show that there exists a non-zero element \(\gamma\) in V such that \((\gamma|\alpha) \geqslant 0\) for all \(\alpha \in P\) . If not, the result of 1) would show that an infinite sequence \(\alpha_{1}, \alpha_{2}, \ldots\) of elements of P could be found such that
\end{enumerate}

\[
\beta_ {i} = \alpha_ {1} + \dots + \alpha_ {i} \in P
\]

for all \(i\) ; there would exist two distinct integers \(i, j\) such that \(\beta_{i} = \beta_{j}\) , which would contradict the result of 2).

\begin{enumerate}
\def\labelenumi{\arabic{enumi})}
\setcounter{enumi}{3}
\tightlist
\item
  To prove the proposition, it is enough (Cor. 2 of Prop. 20) to show that there exists a basis \((\alpha_{k})_{1\leqslant k\leqslant l}\) of \(V\) such that, for the lexicographic order defined by this basis, every element of \(P\) is \(>0\) . We proceed by induction on \(l = \dim V\) , and assume that the proposition is established for all dimensions \(< l\) . Let \(\gamma \in V\) be such that \(\gamma \neq 0\) and \((\gamma |\alpha)\geqslant 0\) for all \(\alpha \in P\) (cf.~3)). Let \(L\) be the hyperplane orthogonal to \(\gamma\) , and \(V'\) the subspace of \(L\) generated by \(R\cap L\) . Then \(R\cap L\) is a root system in \(V'\) and \(P\cap L\) is closed in \(R\cap L\) . By the induction hypothesis, there exists a basis \((\beta_{1},\ldots ,\beta_{l'})\) of \(V'\) such that the elements of \(P\cap L\) are \(>0\) for the lexicographic order defined by this basis. Then any basis of \(V\) whose first \(l' + 1\) elements are \(\gamma ,\beta_{1},\ldots ,\beta_{l'}\) and whose remaining elements are in \(L\) has the required property.
\end{enumerate}

PROPOSITION 23. Let P be a subset of R and \(V_{1}\) (resp. \(\Gamma\) ) the vector subspace (resp. the subgroup) of V generated by P. The following conditions are equivalent:

\begin{enumerate}
\def\labelenumi{(\roman{enumi})}
\item
  P is closed and symmetric;
\item
  P is closed, and P is a root system in \(V_{1}\) ;
\item
  \(\Gamma \cap \mathbb{R} = \mathbb{P}\) .
\end{enumerate}

Assume that these conditions are satisfied. For any \(\alpha \in \mathbb{P}\) , let \(\alpha_{1}^{-}\) be the restriction of \(\alpha^{-}\) to \(\mathrm{V}_{1}\) . Then the map \(\alpha \mapsto \alpha_{1}^{-}\) is the canonical bijection from the root system \(\mathbb{P}\) to \(\mathbb{P}^{-}\) .

\begin{enumerate}
\def\labelenumi{(\roman{enumi})}
\setcounter{enumi}{2}
\item
  \(\Longrightarrow\) (i): clear.
\item
  \(\Longrightarrow\) (ii): assume that P is closed and symmetric. First, P satisfies \((\mathrm{RS}_{\mathrm{I}})\) in \(V_{1}\) . We show that, if \(\alpha, \beta \in P\) , then \(s_{\alpha}(\beta) \in P\) . This is clear if \(\alpha\) and \(\beta\) are proportional. Otherwise, \(s_{\alpha}(\beta) = \beta - n(\beta, \alpha)\alpha\) and \(\beta - p\alpha \in R\) for all rational integers p between 0 and \(n(\beta, \alpha)\) (Prop. 9, no. 3), so
\end{enumerate}

\[
\beta - n (\beta , \alpha) \alpha \in P
\]

since P is closed and symmetric. Thus, \(s_{\alpha,\alpha_{1}}(P) = P\) , and P satisfies \((RS_{II})\) . It is clear that P satisfies \((RS_{III})\) . Thus, P satisfies (ii), and we have proved the last assertion of the proposition at the same time.

\begin{enumerate}
\def\labelenumi{(\roman{enumi})}
\setcounter{enumi}{1}
\tightlist
\item
  \(\Longrightarrow\) (iii): we show that, if condition (ii) is satisfied, then \(\Gamma\cap R=P\) . It is clear that \(P\subseteq\Gamma\cap R\) . Let \(\beta\in\Gamma\cap R\) . Since \(\beta\in\Gamma\) and \(P=-P\) , \(\beta=\alpha_{1}+\alpha_{2}+\cdots+\alpha_{k}\) with \(\alpha_{1},\ldots,\alpha_{k}\in P\) . We shall prove that \(\beta\in P\) . This is clear if k=1. We argue by induction on k. We have
\end{enumerate}

\[
0 <   (\beta | \beta) = \sum_ {i = 1} ^ {k} (\beta | \alpha_ {i}),
\]

so \((\beta |\alpha_{i}) > 0\) for some index \(i\) . If \(\beta = \alpha_{i}\) , then \(\beta \in \mathbb{P}\) . Otherwise, \(\beta - \alpha_{i} \in \mathbb{R}\) (Cor. of Th. 1, no. 3), so \(\beta - \alpha_{i} \in \mathbb{P}\) by the induction hypothesis, hence \(\beta \in \mathbb{P}\) since \(\mathbb{P}\) is closed.

The conditions of Prop. 23 can be realised with \(V_{1} = V\) and yet \(P \neq R\) . For example, this is the case when \(R\) is a system of type \(G_{2}\) and \(P\) a system of type \(A_{2}\) ; cf.~Plate X.

PROPOSITION 24. Let R' be the intersection of R with a vector subspace of V, so that R' is a root system in the vector subspace V' that it generates (cf.~Cor. of Prop. 4, no. 1). Let B' be a basis of R'.

\begin{enumerate}
\def\labelenumi{(\roman{enumi})}
\item
  There exists a basis of R containing \(\mathbf{B}'\) .
\item
  \(\mathbf{R}'\) is the set of elements of \(\mathbf{R}\) that are linear combinations of elements of \(\mathbf{B}'\) .
\end{enumerate}

Assertion (ii) is clear. We prove (i). Let \((\varepsilon_{1},\varepsilon_{2},\ldots,\varepsilon_{l})\) be a basis of V such that \(B' = (\varepsilon_{p+1},\varepsilon_{p+2},\ldots,\varepsilon_{l})\) . The lexicographic order on V corresponding to this basis defines a chamber C of R. It is clear that every element of \(B'\) is minimal in \(\mathrm{R}_{+}(\mathrm{C})\) . Thus \(\mathrm{B}' \subseteq \mathrm{B}(\mathrm{C})\) .

